\begin{proof}[Proof of \cref{obj:lem:exact-calibrations}]
We use the maps in \cref{obj:def:calibration-maps}, with the root-selection and fallback conventions in the statement. For \(r_{\mathrm{amp}}\ge0\), write
\[
\begin{aligned}
\mathcal P_{\mathrm{mix}}(r_{\mathrm{amp}})
&=[-r_{\mathrm{amp}},r_{\mathrm{amp}}]^2\times[0,1],\\
\mathcal P_{\mathrm{fair}}(r_{\mathrm{amp}})
&=[0,1/4]\times[-r_{\mathrm{amp}},r_{\mathrm{amp}}]\times[0,1].
\end{aligned}
\]
These parameter regions are compact. We write \(\mathrm D^m f\) for the \(m\)th total derivative of a parameter function \(f\), with \(\mathrm D^0f=f\).

\begin{enumerate}
\item \emph{Uniform existence and uniqueness of the mixed root.}
Define the mixed residual, wherever its branch denominators are positive, by
\[
\mathcal F_{\mathrm{mix}}(\eta,\zeta,u,\chi)
=
16E\mathfrak B
\bigl(32\eta\zeta,\tfrac12-\eta\chi Z_u,
                    \tfrac12+\zeta Z_u\bigr)-\chi .
\]
The integral defining \(D\) is smooth, and the fundamental theorem of calculus gives
\[
tD(t)=e^t-1,\qquad D(0)=1.
\]
At zero amplitudes the branch radicand is \(1\), its denominator is \(2\), and
\[
\mathfrak B(0,\tfrac12,\tfrac12)=\frac1{16}.
\]
Consequently, for every \(u,\chi\in\mathbb R\),
\[
\mathcal F_{\mathrm{mix}}(0,0,u,\chi)=1-\chi,
\qquad
\partial_\chi\mathcal F_{\mathrm{mix}}(0,0,u,\chi)=-1.
\]
The radicands and denominators remain positive near these points, so the residual is smooth there.

Compactness of \([0,1]\times[7/8,9/8]\), continuity, and the finite number of sign pairs give a positive signed amplitude radius on which all branch denominators are positive and
\[
\left|
\mathcal F_{\mathrm{mix}}(\eta,\zeta,u,\chi)-(1-\chi)
\right|<\frac1{16}
\quad
\left(u\in[0,1],\ \chi\in[7/8,9/8]\right).
\]
In particular,
\[
\mathcal F_{\mathrm{mix}}(\eta,\zeta,u,7/8)>\frac1{16},
\qquad
\mathcal F_{\mathrm{mix}}(\eta,\zeta,u,9/8)<-\frac1{16}.
\]
The intermediate value theorem supplies a zero in \([7/8,9/8]\), hence in \((3/4,5/4)\).

Separately, continuity of the root partial derivative and compactness of
\([0,1]\times[3/4,5/4]\) give another positive signed amplitude radius on which
\[
\partial_\chi\mathcal F_{\mathrm{mix}}(\eta,\zeta,u,\chi)<-\frac12
\quad
\left(u\in[0,1],\ \chi\in[3/4,5/4]\right).
\]
Taking the minimum of these two radii gives
\(\varepsilon_{\mathrm{root}}>0\) such that
\[
\forall(\eta,\zeta,u)\in
\mathcal P_{\mathrm{mix}}(\varepsilon_{\mathrm{root}}),
\qquad
\exists!\chi\in(3/4,5/4):
\mathcal F_{\mathrm{mix}}(\eta,\zeta,u,\chi)=0.
\]
Indeed, the derivative inequality makes the residual strictly decreasing throughout the selection bracket.
% lean: CausalSmith.Stat.LogoddsLowsmoothFrontier.mixedEquation_uniform_existsUnique

\item \emph{Smoothness and derivative bounds for the mixed root and cells.}
At each \((0,0,u)\), the root is \(1\) and its root partial derivative is \(-1\). Apply \cref{obj:lem:scalar-implicit-function} on an open neighborhood where the residual is smooth. The resulting local graph stays in \((3/4,5/4)\) after shrinking its parameter neighborhood. The negative root derivative persists throughout that bracket nearby, so the selected root agrees with this graph.

The same graph is smooth to every finite order. To see the regularity upgrade explicitly, implicit differentiation gives
\[
\mathrm D\mathfrak q(\eta,\zeta,u)
=
-\left.
\frac{
\bigl(\partial_\eta\mathcal F_{\mathrm{mix}},
      \partial_\zeta\mathcal F_{\mathrm{mix}},
      \partial_u\mathcal F_{\mathrm{mix}}\bigr)}
{\partial_\chi\mathcal F_{\mathrm{mix}}}
\right|_{\chi=\mathfrak q(\eta,\zeta,u)} .
\]
On a smaller neighborhood the denominator remains nonzero. If the graph is \(C^k\), this formula makes its first derivative \(C^k\), and therefore makes the graph \(C^{k+1}\). Starting from the \(C^1\) graph supplied by the cited result proves smoothness to every finite order. Compactness of the spatial interval then gives
\(\varepsilon_{\mathrm q}>0\) such that \(\mathfrak q\) is smooth on a neighborhood of every point of
\(\mathcal P_{\mathrm{mix}}(\varepsilon_{\mathrm q})\).

For clarity, define the fully substituted mixed cells. For a family index
\(\iota\in\{0,1\}\), a sign pair
\(\boldsymbol{s}=(s_1,s_2)\in\{-1,1\}^2\), and
\(a,y\in\{0,1\}\), put
\[
\begin{aligned}
e_{\iota}^{\mathrm{loc}}(\eta,\zeta,u,\boldsymbol{s})
&=\frac12+(1-2\iota)\eta\mathfrak q(\eta,\zeta,u)Z_u,\\
m^{\mathrm{loc}}(\zeta,u,\boldsymbol{s})
&=\frac12+\zeta Z_u,\\
c_{\iota}^{\mathrm{loc}}(\eta,\zeta,u,\boldsymbol{s})
&=
\begin{cases}
0,&\iota=0,\\
\mathfrak c(32\eta\zeta,e_1^{\mathrm{loc}},m^{\mathrm{loc}}),
&\iota=1,
\end{cases}\\
\bigl(
\mathcal K^{\mathrm{mix}}_{\iota,\boldsymbol{s},1,1},
\mathcal K^{\mathrm{mix}}_{\iota,\boldsymbol{s},1,0},
\mathcal K^{\mathrm{mix}}_{\iota,\boldsymbol{s},0,1},
\mathcal K^{\mathrm{mix}}_{\iota,\boldsymbol{s},0,0}
\bigr)
&=
\bigl(
e_{\iota}^{\mathrm{loc}}m^{\mathrm{loc}}+c_{\iota}^{\mathrm{loc}},
e_{\iota}^{\mathrm{loc}}(1-m^{\mathrm{loc}})-c_{\iota}^{\mathrm{loc}},
(1-e_{\iota}^{\mathrm{loc}})m^{\mathrm{loc}}-c_{\iota}^{\mathrm{loc}},
(1-e_{\iota}^{\mathrm{loc}})(1-m^{\mathrm{loc}})
+c_{\iota}^{\mathrm{loc}}
\bigr).
\end{aligned}
\]
All arguments on the last line are the displayed parameters and sign pair. At zero amplitudes, the covariance formula is evaluated at
\((0,1/2,1/2)\), where its radicand and denominator are positive. The chain rule therefore makes every displayed cell smooth near each \((0,0,u)\). Compactness in \(u\) and finiteness of the indices give
\(\varepsilon_{\mathrm{cell}}>0\) on whose closed mixed region all these cells are smooth on pointwise neighborhoods.

For each fixed \(m\ge0\), continuity of the total derivatives on the corresponding compact regions gives positive constants satisfying
\[
\|\mathrm D^m\mathfrak q\|\le B_{\mathrm q,m}
\quad\hbox{on }\mathcal P_{\mathrm{mix}}(\varepsilon_{\mathrm q}),
\qquad
\|\mathrm D^m\mathcal K^{\mathrm{mix}}_j\|
\le C^{\mathrm{mix}}_{j,m}
\quad\hbox{on }\mathcal P_{\mathrm{mix}}(\varepsilon_{\mathrm{cell}}).
\]
Here the finite index set is
\[
\mathcal I_{\mathrm{aux}}
=\{0,1\}\times\{-1,1\}^2\times\{0,1\}^2,
\qquad
j=(\iota,\boldsymbol{s},a,y).
\]
Thus the positive constant
\[
B_{\mathrm{cell},m}
=\sum_{j\in\mathcal I_{\mathrm{aux}}}C^{\mathrm{mix}}_{j,m}
\]
bounds the \(m\)th derivative of every mixed cell uniformly over its indices.
% lean: CausalSmith.Stat.LogoddsLowsmoothFrontier.mixedRoot_uniform_smooth_derivative_bounds
% lean: CausalSmith.Stat.LogoddsLowsmoothFrontier.localMixedCell_uniform_smooth_derivative_bounds

\item \emph{The smooth fair residual.}
In the fair calculations use the baseline risk and its normalization
\[
p_*=\frac25,\qquad B_*=1+dp_*,
\qquad d=e^t-1,
\]
and define the unnormalized residual by
\[
\mathcal R_{\mathrm{fair}}(t,\delta,\xi,u)
=
E\mathsf G(t,\xi+\delta Z_u)
-\mathsf G(\mathsf T(t,\delta),\xi).
\]
We first check the domains needed for its derivatives.

For \(0\le t\le1/4\), the elementary exponential bounds
\[
1+t\le e^t\le\frac1{1-t}
\]
give
\[
t\le d\le\frac43t,\qquad
0\le d\le\frac13,\qquad
1\le B_*\le\frac43.
\]
If \(|\delta|\le1/100\), then
\[
\delta^2\le\frac1{10000},
\qquad
d\delta^2\le\frac1{30000},
\]
so both logarithm arguments in \(\mathsf T\) are strictly positive.

The following calculation also supplies the comparator bounds used to justify smooth composition. Define
\[
\mathcal A_{\mathrm{gap}}=\frac{d\delta^2}{B_*}.
\]
The preceding inequalities imply
\[
0\le\mathcal A_{\mathrm{gap}}\le\frac1{1000},
\qquad
\frac34t\delta^2\le\mathcal A_{\mathrm{gap}}
\le\frac43t\delta^2.
\]
The comparator formula becomes
\[
t-\mathsf T(t,\delta)
=
-\log\!\left(1-\frac52\mathcal A_{\mathrm{gap}}\right)
+\log\!\left(1+\frac53\mathcal A_{\mathrm{gap}}\right).
\]
For \(0\le x_{\log},y_{\log}\le1/100\), the inequalities
\(\log z\le z-1\) and \(1-z^{-1}\le\log z\), valid for \(z>0\), yield
\[
x_{\log}+\frac{100}{101}y_{\log}
\le-\log(1-x_{\log})+\log(1+y_{\log})
\le\frac{100}{99}x_{\log}+y_{\log}.
\]
Indeed,
\[
x_{\log}\le-\log(1-x_{\log})
\le\frac{x_{\log}}{1-x_{\log}}
\le\frac{100}{99}x_{\log},
\]
and
\[
\frac{100}{101}y_{\log}
\le\frac{y_{\log}}{1+y_{\log}}
\le\log(1+y_{\log})\le y_{\log}.
\]
Apply these inequalities with
\(x_{\log}=\frac52\mathcal A_{\mathrm{gap}}\) and
\(y_{\log}=\frac53\mathcal A_{\mathrm{gap}}\). Since
\[
\frac34\left(\frac52+\frac{100}{101}\frac53\right)>3,
\qquad
\frac43\left(\frac{100}{99}\frac52+\frac53\right)<6,
\]
we obtain, including \(t=0\) and \(\delta=0\),
\[
3t\delta^2\le t-\mathsf T(t,\delta)\le6t\delta^2,
\qquad
\frac t2\le\mathsf T(t,\delta)\le t.
\]
The last pair follows from \(6\delta^2\le1/2\).

For every real \(u\), the two sign terms give \(|Z_u|\le2\). Thus, for
\(\xi\in[3/10,1/2]\) and \(|\delta|\le1/100\),
\[
|\delta Z_u|\le\frac1{50},
\qquad
\xi+\delta Z_u\ge\frac3{10}-\frac1{50}>0.
\]
Together with \(t\ge0\) and \(\mathsf T(t,\delta)\ge0\), this makes the risk-shift denominators positive. All these strict positivity conditions persist on ambient neighborhoods. Hence \(\mathcal R_{\mathrm{fair}}\) is jointly smooth near every point of the indicated parameter range.

Sign symmetry and the dependence of \(\mathsf T\) on \(\delta^2\) give
\[
\begin{gathered}
\mathcal R_{\mathrm{fair}}(t,-\delta,\xi,u)
=\mathcal R_{\mathrm{fair}}(t,\delta,\xi,u),\\
\mathcal R_{\mathrm{fair}}(t,0,\xi,u)=0,
\qquad
\partial_\delta\mathcal R_{\mathrm{fair}}(t,0,\xi,u)=0,
\qquad
\mathcal R_{\mathrm{fair}}(0,\delta,\xi,u)=0.
\end{gathered}
\]
The last equality uses \(\mathsf T(0,\delta)=0\),
\(\mathsf G(0,\xi)=\xi\), and \(EZ_u=0\).

The normalized extension is the double Taylor integral
\[
\mathfrak H(t,\delta,\xi,u)
=
\int_0^1\int_0^1
(1-\omega)\,
\partial_t\partial_\delta^2
\mathcal R_{\mathrm{fair}}(\tau t,\omega\delta,\xi,u)
\,d\omega\,d\tau.
\]
Here \(\tau,\omega\in[0,1]\) are integration variables. Their compact range permits one ambient neighborhood on which the integrand is smooth for all integration points. Differentiation under the two integrals then proves joint smoothness of \(\mathfrak H\) on that neighborhood, at every finite order.

The fundamental theorem of calculus, using
\(\partial_\delta^2\mathcal R_{\mathrm{fair}}(0,\delta,\xi,u)=0\), gives
\[
t\int_0^1
\partial_t\partial_\delta^2
\mathcal R_{\mathrm{fair}}(\tau t,\omega\delta,\xi,u)\,d\tau
=
\partial_\delta^2
\mathcal R_{\mathrm{fair}}(t,\omega\delta,\xi,u).
\]
The second-order integral Taylor formula gives
\[
\delta^2\int_0^1
(1-\omega)\,
\partial_\delta^2
\mathcal R_{\mathrm{fair}}(t,\omega\delta,\xi,u)\,d\omega
=
\mathcal R_{\mathrm{fair}}(t,\delta,\xi,u).
\]
Interchanging the two integrals therefore proves the undivided identity
\[
t\delta^2\mathfrak H(t,\delta,\xi,u)
=\mathcal R_{\mathrm{fair}}(t,\delta,\xi,u)
\]
throughout
\(t\in[0,1/4]\), \(|\delta|\le1/100\),
\(\xi\in[3/10,1/2]\), and \(u\in[0,1]\).
For \(t\delta\ne0\), division gives the quotient in the statement.
% lean: CausalSmith.Stat.LogoddsLowsmoothFrontier.comparatorEffect_gap_bounds
% lean: CausalSmith.Stat.LogoddsLowsmoothFrontier.comparatorEffect_range_bounds
% lean: CausalSmith.Stat.LogoddsLowsmoothFrontier.fairEquation_mul_parameters
% lean: CausalSmith.Stat.LogoddsLowsmoothFrontier.fairEquation_eq_div

\item \emph{Fair endpoint factorization and the removable axes.}
At either spatial endpoint the law of \(Z_u\) is one fair sign. Exponentiating the comparator formula gives
\[
e^{\mathsf T(t,\delta)}
=e^t
\frac{1-d\delta^2/(p_*B_*)}
     {1+d\delta^2/((1-p_*)B_*)},
\]
or, equivalently,
\[
e^{\mathsf T(t,\delta)}
\left(\frac35B_*+d\delta^2\right)
=
e^t\left(\frac35B_*-\frac32d\delta^2\right).
\]
Clearing the positive risk denominators using this identity proves
\[
\frac{\mathsf G(t,p_*+\delta)+\mathsf G(t,p_*-\delta)}2
=\mathsf G(\mathsf T(t,\delta),p_*).
\]
It follows first for \(t\delta\ne0\) that
\(\mathfrak H(t,\delta,p_*,u)=0\) at \(u=0,1\).
For each nonzero \(\delta\), continuity in \(t\) extends this equality to \(t=0\); continuity in \(\delta\) then extends it to \(\delta=0\). Thus
\[
\mathfrak H(t,\delta,p_*,0)
=\mathfrak H(t,\delta,p_*,1)=0
\quad
\left(t\in[0,1/4],\ |\delta|\le1/100\right).
\]

To identify every endpoint zero, define
\[
\begin{aligned}
\mathcal S_{\mathrm{end}}(d,\delta,\xi)
&=-25d\delta^2+25d\xi^2-15d\xi-6d+25\xi-15,\\
\mathcal J_{\mathrm{end}}(t,\delta,\xi)
&=\bigl(1+d(\xi+\delta)\bigr)
  \bigl(1+d(\xi-\delta)\bigr)
  \bigl(1+(e^{\mathsf T(t,\delta)}-1)\xi\bigr)
  \left(\frac35B_*+d\delta^2\right).
\end{aligned}
\]
The second expression is positive on the fair range
\(t\in[0,1/4]\), \(|\delta|\le1/100\),
\(\xi\in[3/10,1/2]\). Directly clearing the same denominators and substituting the corrected exponential identity above gives
\[
\mathcal R_{\mathrm{fair}}(t,\delta,\xi,0)
\mathcal J_{\mathrm{end}}(t,\delta,\xi)
=
-\frac{d\delta^2e^t}{50}
(5\xi-2)\mathcal S_{\mathrm{end}}(d,\delta,\xi).
\]
For \(\delta\ne0\), the undivided Taylor identity cancels the amplitude square:
\[
t\mathfrak H(t,\delta,\xi,0)
\mathcal J_{\mathrm{end}}(t,\delta,\xi)
=
-\frac{de^t}{50}
(5\xi-2)\mathcal S_{\mathrm{end}}(d,\delta,\xi).
\]
Continuity in \(\delta\) extends this identity to zero amplitude, where it reads
\[
t\mathfrak H(t,0,\xi,0)(1+d\xi)^3\frac35B_*
=
-\frac{de^t}{50}
(5\xi-2)\mathcal S_{\mathrm{end}}(d,0,\xi).
\]
For \(d\ge0\) and \(\xi\in[3/10,1/2]\), the inequality
\(\xi^2\le\xi/2\) gives
\[
\mathcal S_{\mathrm{end}}(d,\delta,\xi)
\le-\frac52d\xi-6d+25\xi-15
\le-\frac52<0.
\]
Thus, when \(t>0\), the zero-amplitude endpoint residual has the sign of \(5\xi-2\). At either endpoint, when also \(\delta\ne0\), the factored unnormalized identity shows that every matching root in the selection bracket equals \(p_*\).

We also need position independence at zero amplitude. The sign moments are
\[
EZ_u=0,\qquad EZ_u^2
=\cos^2(\pi u/2)+\sin^2(\pi u/2)=1.
\]
Differentiating the averaged risk twice in amplitude gives
\[
\left.
\partial_\delta^2 E\mathsf G(t,\xi+\delta Z_u)
\right|_{\delta=0}
=
-\frac{2e^td}{(1+d\xi)^3}.
\]
The comparator term has no spatial argument, so
\[
\partial_\delta^2
\mathcal R_{\mathrm{fair}}(t,0,\xi,u)
=
\partial_\delta^2
\mathcal R_{\mathrm{fair}}(t,0,\xi,0).
\]
Their effect derivatives agree as well, including at the endpoints of the effect interval by smoothness and the one-sided derivative there. Substitution in the double integral proves
\[
\mathfrak H(t,0,\xi,u)=\mathfrak H(t,0,\xi,0).
\]

For zero effect, differentiation of the explicit comparator gives
\[
\partial_t\mathsf T(0,\delta)=1-\frac{\delta^2}{6/25},
\qquad
\partial_t\mathsf G(0,\xi)=\xi(1-\xi).
\]
Using the two sign moments,
\[
\partial_t\mathcal R_{\mathrm{fair}}(0,\delta,\xi,u)
=
\delta^2\left(-1+\frac{\xi(1-\xi)}{6/25}\right).
\]
Differentiate the undivided Taylor identity at \(t=0\). For \(\delta\ne0\), cancellation of \(\delta^2\) gives the first identity below; continuity in \(\delta\) supplies the same identity at \(\delta=0\):
\[
\mathfrak H(0,\delta,\xi,u)
=-1+\frac{\xi(1-\xi)}{6/25},
\qquad
\mathfrak H(0,\delta,\xi,u)=0
\ \Longleftrightarrow\
(\xi-\tfrac25)(\xi-\tfrac35)=0.
\]
Within \((3/10,1/2)\), the zero is therefore \(p_*\).
Combining these facts yields
\[
\mathfrak H(t,0,p_*,u)=0,
\qquad
\mathfrak H(t,0,\xi,u)=0,\quad \xi\in(3/10,1/2)
\ \Longrightarrow\ \xi=p_*.
\]
Here \(t=0\) uses the explicit quadratic formula, and \(t>0\) uses the negative residual factor.
% lean: CausalSmith.Stat.LogoddsLowsmoothFrontier.fairEquation_endpoint_factor_zero_amplitude
% lean: CausalSmith.Stat.LogoddsLowsmoothFrontier.fairEquation_zero_amplitude_unique
% lean: CausalSmith.Stat.LogoddsLowsmoothFrontier.fairEquation_zero_effect

\item \emph{Uniform fair-root existence.}
The preceding factorization shows opposite bracket signs at zero amplitude when \(t>0\). At \(t=0\), the explicit extension gives
\[
\mathfrak H(0,0,3/10,u)=-\frac18,
\qquad
\mathfrak H(0,0,1/2,u)=\frac1{24}.
\]
Hence throughout the compact effect and spatial ranges,
\[
\mathfrak H(t,0,3/10,u)<0,
\qquad
\mathfrak H(t,0,1/2,u)>0.
\]
Joint continuity and compactness preserve these strict inequalities for
\(|\delta|\le\varepsilon_{\mathrm{exist}}\), for some
\[
0<\varepsilon_{\mathrm{exist}}\le\frac1{100}.
\]
Continuity in the margin and the intermediate value theorem then give
\[
\forall(t,\delta,u)\in
\mathcal P_{\mathrm{fair}}(\varepsilon_{\mathrm{exist}}),
\qquad
\exists\xi\in(3/10,1/2):
\mathfrak H(t,\delta,\xi,u)=0.
\]
The strict endpoint signs put the zero in the open bracket.
% lean: CausalSmith.Stat.LogoddsLowsmoothFrontier.fairEquation_uniform_exists

\item \emph{Smoothness of the selected fair root.}
Differentiate the zero-amplitude factorization at \(\xi=p_*\), where the residual itself vanishes. This gives
\[
t\,\partial_\xi\mathfrak H(t,0,p_*,u)
B_*^3\frac35B_*
=
\frac{de^t(8d+5)}{10}.
\]
For \(t>0\), all factors except the derivative on the left are positive, and the right side is positive. At zero effect, the explicit extension gives
\[
\partial_\xi\mathfrak H(0,0,p_*,u)=\frac56.
\]
Therefore the center slope is positive throughout the full effect and spatial ranges.

Continuity of this partial derivative and compactness give a radius
\(r_{\mathrm{center}}\) satisfying
\[
0<r_{\mathrm{center}}\le\frac1{100},
\qquad
\partial_\xi\mathfrak H(t,\delta,\xi,u)>0
\]
whenever
\[
t\in[0,1/4],\quad
|\delta|\le r_{\mathrm{center}},\quad
|\xi-p_*|\le r_{\mathrm{center}},\quad
u\in[0,1].
\]
The zero-amplitude uniqueness from the preceding steps also confines all nearby roots. Indeed, the set
\[
\left\{
(t,u,\xi):
t\in[0,1/4],\ u\in[0,1],\
\xi\in[3/10,1/2],\
|\xi-p_*|\ge r_{\mathrm{center}}
\right\}
\]
is compact, and \(\mathfrak H(t,0,\xi,u)\) is nonzero on it. Continuity therefore gives
\(0<\varepsilon_{\mathrm{conf}}\le1/100\) such that
\[
\mathfrak H(t,\delta,\xi,u)=0,\quad
\xi\in[3/10,1/2],\quad
|\delta|\le\varepsilon_{\mathrm{conf}}
\quad\Longrightarrow\quad
|\xi-p_*|<r_{\mathrm{center}}.
\]
If the displayed compact set is empty, the implication follows directly from its defining condition.

Take the minimum of
\(r_{\mathrm{center}},\varepsilon_{\mathrm{conf}}\), and
\(\varepsilon_{\mathrm{exist}}\). Every selected root on the resulting closed fair region is a genuine bracket zero, lies in the center neighborhood, and has nonzero margin derivative. Moreover, these controls give uniqueness on ambient parameter neighborhoods: positivity persists on the compact center bracket, and nonvanishing persists on the compact portion of the larger bracket outside it. Any two nearby bracket zeros must therefore lie in the center bracket, where the residual is strictly increasing.

Apply \cref{obj:lem:scalar-implicit-function} at each selected zero. Nearby bracket uniqueness identifies the selected root with the local graph. The first-derivative identity is
\[
\mathrm D\mathfrak p(t,\delta,u)
=
-\left.
\frac{
\bigl(\partial_t\mathfrak H,
      \partial_\delta\mathfrak H,
      \partial_u\mathfrak H\bigr)}
{\partial_\xi\mathfrak H}
\right|_{\xi=\mathfrak p(t,\delta,u)} .
\]
The regularity upgrade used for the mixed graph applies again. Thus there is
\[
0<\varepsilon_{\mathrm{smooth}}\le\frac1{100}
\]
such that \(\mathfrak p\) is smooth on a neighborhood of every point of
\(\mathcal P_{\mathrm{fair}}(\varepsilon_{\mathrm{smooth}})\).
% lean: CausalSmith.Stat.LogoddsLowsmoothFrontier.fairRoot_uniform_contDiffAt

\item \emph{A preliminary common radius and the quadratic fair bound.}
Define
\[
\varepsilon_{\mathrm{pre}}
=
\min\{\varepsilon_{\mathrm{smooth}},
      \varepsilon_{\mathrm q},
      \varepsilon_{\mathrm{cell}}\},
\qquad
\varepsilon_0
=
\min\{\varepsilon_{\mathrm{pre}},
      \varepsilon_{\mathrm{exist}}\}.
\]
Then
\[
0<\varepsilon_0\le\frac1{100},
\qquad
\varepsilon_0\le\varepsilon_{\mathrm q},
\qquad
\varepsilon_0\le\varepsilon_{\mathrm{cell}},
\]
and \(\mathfrak p\) is smooth on pointwise neighborhoods of its closed region.

The endpoint and axis computations identify the selected roots there:
\[
\mathfrak p(t,\delta,0)=p_*,
\qquad
\mathfrak p(t,0,u)=p_*,
\qquad
\mathfrak p(0,\delta,u)=p_*.
\]
For the first equality, split first on \(t=0\), then on \(\delta=0\); the zero-effect and zero-amplitude uniqueness computations settle these cases. When both are nonzero, endpoint matching and the negative residual factor force the selected bracket root to be \(p_*\).

The numerator is even in \(\delta\). Its second amplitude derivative is consequently even, so its effect derivative and the double integral satisfy
\[
\mathfrak H(t,-\delta,\xi,u)=\mathfrak H(t,\delta,\xi,u).
\]
The selectors for these identical equations have the same value, including their common fallback convention. Hence
\[
\mathfrak p(t,-\delta,u)=\mathfrak p(t,\delta,u).
\]

Smoothness and compactness give positive bounds for the two partial derivatives needed below. More explicitly, define
\[
\begin{aligned}
C_{\delta\delta}
&=
\max_{\mathcal P_{\mathrm{fair}}(\varepsilon_0)}
|\partial_\delta^2\mathfrak p|,\\
C_{\delta\delta u}
&=
\max_{\mathcal P_{\mathrm{fair}}(\varepsilon_0)}
|\partial_\delta^2\partial_u\mathfrak p|,\\
B_{\mathrm{quad}}
&=\max\{1,C_{\delta\delta},C_{\delta\delta u}\},
\qquad M=2B_{\mathrm{quad}}>0.
\end{aligned}
\]
Both maxima exist because the indicated partial derivatives are continuous on the compact region.

For a smooth even function \(\varphi\), its derivative at zero vanishes, and the integral Taylor formula gives
\[
\varphi(\delta)-\varphi(0)
=
\delta^2\int_0^1(1-\omega)\varphi''(\omega\delta)\,d\omega.
\]
If \(|\varphi''|\le B_{\mathrm{quad}}\) on the segment joining \(0\) to \(\delta\), then
\[
|\varphi(\delta)-\varphi(0)|
\le B_{\mathrm{quad}}\delta^2.
\]
Every amplitude on that segment has absolute value at most
\(|\delta|\le\varepsilon_0\).

Apply this first to
\(\varphi(\delta')=\mathfrak p(t,\delta',u)\), and then to
\(\varphi(\delta')=\partial_u\mathfrak p(t,\delta',u)\), with \(t,u\) fixed in their prescribed ranges. Both functions are smooth and even. Their zero-amplitude values are respectively \(p_*\) and \(0\). For the latter value, constancy of
\(\mathfrak p(t,0,u)=p_*\) on \([0,1]\) forces the spatial derivative to vanish also at \(u=0,1\): the available two-sided derivative equals the zero one-sided derivative from the interval. Therefore
\[
|\mathfrak p(t,\delta,u)-p_*|
\le B_{\mathrm{quad}}\delta^2,
\qquad
|\partial_u\mathfrak p(t,\delta,u)|
\le B_{\mathrm{quad}}\delta^2,
\]
and hence
\[
|\mathfrak p(t,\delta,u)-p_*|
+|\partial_u\mathfrak p(t,\delta,u)|
\le M\delta^2
\quad\hbox{on }\mathcal P_{\mathrm{fair}}(\varepsilon_0).
\]
% lean: CausalSmith.Stat.LogoddsLowsmoothFrontier.fairRoot_endpoint_center
% lean: CausalSmith.Stat.LogoddsLowsmoothFrontier.fairRoot_quadratic_bound_of_smooth

\item \emph{Fair-cell regularity and common derivative bounds.}
For the same auxiliary indices as before, define
\[
\begin{aligned}
\mu^{\mathrm{loc}}_{0,\iota}
&=
\begin{cases}
\mathfrak p(t,\delta,u),&\iota=0,\\
\mathfrak p(t,\delta,u)+\delta Z_u,&\iota=1,
\end{cases}\\
t^{\mathrm{loc}}_\iota
&=
\begin{cases}
\mathsf T(t,\delta),&\iota=0,\\
t,&\iota=1,
\end{cases}\\
\mu^{\mathrm{loc}}_{1,\iota}
&=\mathsf G(t^{\mathrm{loc}}_\iota,\mu^{\mathrm{loc}}_{0,\iota}),\\
\bigl(
\mathcal K^{\mathrm{fair}}_{\iota,\boldsymbol{s},1,1},
\mathcal K^{\mathrm{fair}}_{\iota,\boldsymbol{s},1,0},
\mathcal K^{\mathrm{fair}}_{\iota,\boldsymbol{s},0,1},
\mathcal K^{\mathrm{fair}}_{\iota,\boldsymbol{s},0,0}
\bigr)
&=
\frac12\bigl(
\mu^{\mathrm{loc}}_{1,\iota},
1-\mu^{\mathrm{loc}}_{1,\iota},
\mu^{\mathrm{loc}}_{0,\iota},
1-\mu^{\mathrm{loc}}_{0,\iota}
\bigr).
\end{aligned}
\]
These are functions of \((t,\delta,u)\), with the signs evaluated in \(Z_u\).

On every input the fair selector belongs to \((3/10,1/2)\): it either selects a root there or takes the fallback value \(p_*\), which is also in that bracket. Consequently, when \(|\delta|\le1/100\),
\[
\frac14\le\mu^{\mathrm{loc}}_{0,\iota}\le\frac35.
\]
Indeed, \(|\delta Z_u|\le1/50\), and adding this bound to either side of the selection bracket proves the displayed bounds. The local effect is nonnegative by the comparator range estimate. Hence
\[
1+\bigl(e^{t^{\mathrm{loc}}_\iota}-1\bigr)
\mu^{\mathrm{loc}}_{0,\iota}>0.
\]
The chain rule, root smoothness, and positivity of the logarithm arguments now prove smoothness of every fully substituted fair cell on neighborhoods of the closed fair region.

For each \(m\ge0\), compactness gives positive bounds
\[
\|\mathrm D^m\mathfrak p\|\le B_{\mathrm p,m},
\qquad
\|\mathrm D^m\mathcal K^{\mathrm{fair}}_j\|
\le C^{\mathrm{fair}}_{j,m}
\quad\hbox{on }\mathcal P_{\mathrm{fair}}(\varepsilon_0).
\]
Set
\[
\begin{aligned}
B_{\mathrm{fair},m}
&=B_{\mathrm p,m}
  +\sum_{j\in\mathcal I_{\mathrm{aux}}}C^{\mathrm{fair}}_{j,m},\\
B_{\mathrm{rem},m}
&=\max\{B_{\mathrm{cell},m},B_{\mathrm{fair},m}\},\\
B_m&=\max\{B_{\mathrm q,m},B_{\mathrm{rem},m}\}>0.
\end{aligned}
\]
Restriction to the smaller mixed region
\(\mathcal P_{\mathrm{mix}}(\varepsilon_0)\) preserves its earlier bounds. Thus
\[
\begin{array}{ll}
\|\mathrm D^m\mathfrak q\|\le B_m,\quad
\|\mathrm D^m\mathcal K^{\mathrm{mix}}_j\|\le B_m
&\text{on }\mathcal P_{\mathrm{mix}}(\varepsilon_0),\\[2mm]
\|\mathrm D^m\mathfrak p\|\le B_m,\quad
\|\mathrm D^m\mathcal K^{\mathrm{fair}}_j\|\le B_m
&\text{on }\mathcal P_{\mathrm{fair}}(\varepsilon_0).
\end{array}
\]
All these functions are \(C^\infty\) on their respective closed regions.
% lean: CausalSmith.Stat.LogoddsLowsmoothFrontier.localFairCell_smooth_derivative_bounds_of_root
% lean: CausalSmith.Stat.LogoddsLowsmoothFrontier.exact_calibrations

\item \emph{Fixed open branch neighborhoods.}
We construct open sets for the actual selectors and cells. Define
\[
\begin{aligned}
\mathcal A_{\mathrm{mix}}
=\bigl\{(\eta,\zeta,u)\in\mathbb R^3:\;&
\mathfrak q(\eta,\zeta,u)
\text{ is the unique zero of }
\mathcal F_{\mathrm{mix}}(\eta,\zeta,u,\cdot)
\text{ in }(3/4,5/4),\\[-1mm]
&\mathfrak q\text{ and every }\mathcal K^{\mathrm{mix}}_j
\text{ are smooth at }(\eta,\zeta,u)\bigr\},\\
U&=\operatorname{int}\mathcal A_{\mathrm{mix}}.
\end{aligned}
\]
At every \((0,0,u)\), the local implicit graph through \(1\), its bracket range, and the persistent negative root derivative give the unique selected zero nearby. The root and all mixed cells are smooth on those same neighborhoods. Therefore
\[
(0,0,u)\in U\qquad(u\in[0,1]).
\]
Openness and compactness of this spatial segment give
\(\varepsilon_{\mathrm{mix}}>0\) such that
\[
\mathcal P_{\mathrm{mix}}(\varepsilon_{\mathrm{mix}})\subset U.
\]
By the defining conditions, the mixed selector is the unique bracket zero throughout \(U\), and it and all mixed cells are smooth there.

For the fair branch fix
\[
b=\frac1{20},\qquad 0<b<\frac1{10},
\]
and define
\[
\begin{aligned}
\mathcal A_{\mathrm{fair}}
=\bigl\{(t,\delta,u)\in\mathbb R^3:\;&
\mathfrak p(t,\delta,u)\in(p_*-b,p_*+b)\cap(3/10,1/2),\\
&\mathfrak p(t,\delta,u)
\text{ is the unique zero of }
\mathfrak H(t,\delta,\cdot,u)
\text{ in }(p_*-b,p_*+b),\\
&\mathfrak p\text{ and every }\mathcal K^{\mathrm{fair}}_j
\text{ are smooth at }(t,\delta,u)\bigr\},\\
V&=\operatorname{int}\mathcal A_{\mathrm{fair}}.
\end{aligned}
\]
At each \((t,0,u)\) in the full effect and spatial ranges, the local implicit graph starts at \(p_*\), and therefore stays in
\((p_*-b,p_*+b)\) nearby. The ambient uniqueness argument from the center derivative and compact root exclusion identifies it with the selector from the larger bracket. Smoothness of the residual holds on an ambient open set; the derivative upgrade applies to this same graph at every nearby point. The outer cell formulas have positive denominators and logarithm arguments at zero amplitude, so all fair cells are smooth nearby as well. Consequently,
\[
(t,0,u)\in V
\qquad(t\in[0,1/4],\ u\in[0,1]).
\]
Compactness gives a positive amplitude radius for which the corresponding closed fair region lies in \(V\). Reduce that radius to obtain
\[
0<\varepsilon_{\mathrm{fair}}\le\frac1{100},
\qquad
\mathcal P_{\mathrm{fair}}(\varepsilon_{\mathrm{fair}})\subset V.
\]

For the residual itself define
\[
W=
\operatorname{int}
\left\{
(t,\delta,\xi,u)\in\mathbb R^4:
\mathfrak H\text{ is smooth at }(t,\delta,\xi,u)
\right\}.
\]
The compact integration ranges in its Taylor representation give a common neighborhood of smoothness around every point of
\[
t\in[0,1/4],\quad |\delta|\le1/100,\quad
\xi\in[3/10,1/2],\quad u\in[0,1].
\]
Since
\((p_*-b,p_*+b)\subset[3/10,1/2]\), the required centered fair core with amplitude radius \(\varepsilon_{\mathrm{fair}}\) is contained in \(W\). The residual is smooth on \(W\), and its quotient identity on this core follows from the undivided Taylor identity.

Finally define
\[
\varepsilon_{\mathrm{branch}}
=\frac12\min\{\varepsilon_{\mathrm{mix}},
             \varepsilon_{\mathrm{fair}}\}>0.
\]
Then
\[
\varepsilon_{\mathrm{branch}}<\varepsilon_{\mathrm{mix}},
\qquad
\varepsilon_{\mathrm{branch}}<\varepsilon_{\mathrm{fair}},
\]
which supplies the strict separation between the eventual closed-region radius and both open-neighborhood radii.
% lean: CausalSmith.Stat.LogoddsLowsmoothFrontier.calibration_branch_neighbourhood_exists

\item \emph{The final common radius.}
Set
\[
\varepsilon
=
\min\left\{
\varepsilon_0,\varepsilon_{\mathrm{root}},
\frac1{100},\varepsilon_{\mathrm{branch}}
\right\}>0.
\]
In particular,
\[
\varepsilon\le\varepsilon_0,\qquad
\varepsilon\le\varepsilon_{\mathrm{root}},\qquad
\varepsilon\le\frac1{100}\le\frac18,
\]
and
\[
\varepsilon<\varepsilon_{\mathrm{mix}},
\qquad
\varepsilon<\varepsilon_{\mathrm{fair}}.
\]
The inclusions
\[
\mathcal P_{\mathrm{mix}}(\varepsilon)
\subset\mathcal P_{\mathrm{mix}}(\varepsilon_0),
\qquad
\mathcal P_{\mathrm{fair}}(\varepsilon)
\subset\mathcal P_{\mathrm{fair}}(\varepsilon_0)
\]
preserve smoothness, every derivative bound \(B_m\), and the quadratic estimate with the same positive \(M\). The open sets and the larger signed radii remain as constructed.
% lean: CausalSmith.Stat.LogoddsLowsmoothFrontier.exact_calibrations

\item \emph{Positive odds tables.}
Suppose
\[
|t|\le\varepsilon,\qquad
|\xi-\tfrac12|\le\varepsilon,\qquad
|\upsilon-\tfrac12|\le\varepsilon.
\]
Since \(\varepsilon\le1/8\), the standard estimate
\(|e^t-1|\le2|t|\) for \(|t|\le1\) gives
\[
|d|\le\frac14,\qquad
\frac38\le\xi,\upsilon,1-\xi,1-\upsilon\le\frac58.
\]
Write the branch radicand as
\[
\Delta_{\mathrm{cov}}=L_*^2-4d^2v_*.
\]
The elementary variance and margin-product bounds give
\[
0\le v_*\le\frac1{16},
\qquad
0\le\xi(1-\upsilon)+(1-\xi)\upsilon\le1,
\]
and therefore
\[
L_*\ge\frac34,\qquad
d^2\le\frac1{16},\qquad
\Delta_{\mathrm{cov}}\ge0,\qquad
L_*+\sqrt{\Delta_{\mathrm{cov}}}\ge\frac34.
\]
The covariance formula consequently satisfies
\[
|\mathfrak c(t,\xi,\upsilon)|
=
\frac{2|d|v_*}{L_*+\sqrt{\Delta_{\mathrm{cov}}}}
\le
\frac{2(1/4)(1/16)}{3/4}
=\frac1{24}.
\]
Every unadjusted cell product is at least \(9/64\). Thus each adjusted cell is at least
\[
\frac9{64}-\frac1{24}>0.
\]
Direct addition gives both prescribed margins and total mass one.

Put \(c=\mathfrak c(t,\xi,\upsilon)\). Squaring the square root and clearing its positive branch denominator gives
\[
dc^2-L_*c+dv_*=0.
\]
Expanding the four cells now yields
\[
p_{11}p_{00}-e^t p_{10}p_{01}
=-(dc^2-L_*c+dv_*)=0.
\]
Since \(p_{10}p_{01}>0\), division proves the stated odds ratio.

Finally, \(d=tD(t)\) in the branch formulas gives
\[
\mathfrak c(t,\xi,\upsilon)=t\mathfrak B(t,\xi,\upsilon).
\]
Thus \(\mathfrak B=\mathfrak c/t\) when \(t\ne0\). At \(t=0\), the branch denominator is \(2\) and \(D(0)=1\), so
\[
\mathfrak B(0,\xi,\upsilon)
=\xi(1-\xi)\upsilon(1-\upsilon).
\]
% lean: CausalSmith.Stat.LogoddsLowsmoothFrontier.covarianceBranch_table_positive
% lean: CausalSmith.Stat.LogoddsLowsmoothFrontier.covarianceBranch_odds_ratio
% lean: CausalSmith.Stat.LogoddsLowsmoothFrontier.normalizedBranch_eq_div
% lean: CausalSmith.Stat.LogoddsLowsmoothFrontier.normalizedBranch_zero

\item \emph{The signed mixed identities.}
For \((\eta,\zeta,u)\in\mathcal P_{\mathrm{mix}}(\varepsilon)\), the established bracket-root existence makes the selector a genuine root:
\[
\frac34<\mathfrak q(\eta,\zeta,u)<\frac54,
\qquad
\mathcal F_{\mathrm{mix}}
(\eta,\zeta,u,\mathfrak q(\eta,\zeta,u))=0.
\]
The two endpoint sign laws give identical residual functions and therefore identical selected values:
\[
\mathfrak q(\eta,\zeta,0)=\mathfrak q(\eta,\zeta,1).
\]
Using \(\mathfrak c=t\mathfrak B\), multiplication of the normalized root equation gives
\[
\begin{aligned}
&E\mathfrak c
\bigl(32\eta\zeta,
\tfrac12-\eta\mathfrak q(\eta,\zeta,u)Z_u,
\tfrac12+\zeta Z_u\bigr)\\
&\hspace{15mm}
=32\eta\zeta\,E\mathfrak B
\bigl(32\eta\zeta,
\tfrac12-\eta\mathfrak q(\eta,\zeta,u)Z_u,
\tfrac12+\zeta Z_u\bigr)
=2\eta\zeta\mathfrak q(\eta,\zeta,u).
\end{aligned}
\]
This multiplication also covers either zero-amplitude axis.

The separate sign symmetries follow algebraically from the normalized branch. First,
\[
D(-t)=\frac{D(t)}{e^t};
\]
for \(t\ne0\) this follows from \(tD(t)=e^t-1\), and at \(t=0\) it follows from \(D(0)=1\). Under
\((t,\xi,\upsilon)\mapsto(-t,1-\xi,\upsilon)\), the branch coefficients transform as
\[
d\mapsto-\frac d{e^t},
\qquad
v_*\mapsto v_*,
\qquad
L_*\mapsto\frac{L_*}{e^t}.
\]
The square root and the whole denominator are consequently divided by the positive number \(e^t\), just as the normalized numerator is. This proves
\[
\mathfrak B(-t,1-\xi,\upsilon)
=\mathfrak B(t,\xi,\upsilon).
\]
Exchanging the two margins gives
\[
\mathfrak B(-t,\xi,1-\upsilon)
=\mathfrak B(t,\xi,\upsilon).
\]
Substitution in the mixed residual gives, with the root argument fixed,
\[
\mathcal F_{\mathrm{mix}}(-\eta,\zeta,u,\chi)
=\mathcal F_{\mathrm{mix}}(\eta,\zeta,u,\chi)
=\mathcal F_{\mathrm{mix}}(\eta,-\zeta,u,\chi).
\]
Identical residuals give identical selectors, so the two asserted root symmetries follow.

On the axes, the zero-effect normalized formula and \(EZ_u^2=1\) reduce the residuals to
\[
\mathcal F_{\mathrm{mix}}(0,\zeta,u,\chi)
=1-4\zeta^2-\chi,
\qquad
\mathcal F_{\mathrm{mix}}(\eta,0,u,\chi)
=1-4\eta^2\chi^2-\chi.
\]
When \(|\zeta|\le1/8\), the candidate \(1-4\zeta^2\) lies strictly in \((3/4,5/4)\); the affine equation forces every selected zero to equal it. Hence
\[
\mathfrak q(0,\zeta,u)=1-4\zeta^2.
\]
For the other axis, \(\eta^2\le1/64\) gives
\[
1-4\eta^2(7/8)^2-\frac78\ge0,
\qquad
1-4\eta^2-1\le0.
\]
The intermediate value theorem supplies a zero in \([7/8,1]\), hence in the selection bracket. Substitution of the selected zero into the reduced equation proves
\[
\mathfrak q(\eta,0,u)
+4\eta^2\mathfrak q(\eta,0,u)^2=1.
\]
These arguments include \(\eta=0\) and \(\zeta=0\).
% lean: CausalSmith.Stat.LogoddsLowsmoothFrontier.mixedRoot_spec
% lean: CausalSmith.Stat.LogoddsLowsmoothFrontier.mixedEquation_covariance_matching
% lean: CausalSmith.Stat.LogoddsLowsmoothFrontier.mixedRoot_even_left
% lean: CausalSmith.Stat.LogoddsLowsmoothFrontier.mixedRoot_even_right
% lean: CausalSmith.Stat.LogoddsLowsmoothFrontier.mixedRoot_zero_left
% lean: CausalSmith.Stat.LogoddsLowsmoothFrontier.mixedRoot_zero_right

\item \emph{The fair identities and separation.}
Fix
\[
t\in[0,1/4],\qquad |\delta|\le\varepsilon,\qquad u\in[0,1],
\]
and put \(p=\mathfrak p(t,\delta,u)\). The comparator calculation already gives
\[
3t\delta^2\le t-\mathsf T(t,\delta)\le6t\delta^2,
\qquad
\frac t2\le\mathsf T(t,\delta)\le t.
\]
The preliminary radius construction supplies a bracket root and the identities
\[
\mathfrak p(t,\delta,0)=p_*,
\qquad
\mathfrak p(t,0,u)=p_*,
\qquad
\mathfrak p(0,\delta,u)=p_*.
\]
The endpoint sign laws give identical fair residuals at \(u=0,1\), hence
\[
\mathfrak p(t,\delta,0)=\mathfrak p(t,\delta,1)=p_*.
\]
Evenness supplies \(\mathfrak p(t,-\delta,u)=p\), and the previously obtained quadratic bound gives
\[
|p-p_*|+|\partial_u\mathfrak p(t,\delta,u)|
\le M\delta^2.
\]

To verify the normalized equation on the entire range, first consider \(\delta=0\). The zero-amplitude center is a genuine bracket zero. If \(\delta\ne0\) and \(u\in\{0,1\}\), the endpoint center is a genuine bracket zero by the continuity argument through both axes. In the remaining case, uniform bracket-root existence and the selector definition give the root equation directly. Thus in all cases
\[
\mathfrak H(t,\delta,p,u)=0.
\]
The selector lies in \((3/10,1/2)\), so the fair Taylor identity applies at this margin. For \(t\delta\ne0\), it gives
\[
\mathfrak H(t,\delta,p,u)
=
\frac{\mathcal R_{\mathrm{fair}}(t,\delta,p,u)}{t\delta^2}.
\]
For all \(t,\delta\), including the axes, its undivided form gives
\[
\mathcal R_{\mathrm{fair}}(t,\delta,p,u)
=t\delta^2\mathfrak H(t,\delta,p,u)=0.
\]
By the definition of the numerator, this is exactly
\[
E\mathsf G(t,p+\delta Z_u)
=\mathsf G(\mathsf T(t,\delta),p).
\]
Together with the closed-region regularity, derivative bounds, and open neighborhoods already constructed, these are all the asserted properties of the positive constants \(\varepsilon\) and \(M\).
% lean: CausalSmith.Stat.LogoddsLowsmoothFrontier.fair_matching_of_equation
% lean: CausalSmith.Stat.LogoddsLowsmoothFrontier.exact_calibrations
\end{enumerate}
\end{proof}
